\documentclass[a4paper,12pt]{article}
% [roman] numbers questions as 1(ii)(c); the default [alph] numbers as
% 1(b)(iii). [pdfbookmark] adds a PDF bookmark per question.
\usepackage[roman,pdfbookmark]{ou-tma}

\myname{Alex Morgan}
\mypin{A1234567}
\mycourse{M208: Pure Mathematics}
\mytma{03}

\begin{document}
\maketitle

\begin{question}
\qpart
Prove, by induction on $n$, that
\[
  \sum_{k=1}^{n} k^2 = \frac{n(n+1)(2n+1)}{6}
\]
for all $n \geq 1$.

\emph{Base case.} For $n=1$, the left side is $1$ and the right side is
$\frac{1\cdot 2\cdot 3}{6} = 1$, so the statement holds.

\emph{Inductive step.} Assume the identity holds for some $n \geq 1$.
Adding the next term,
\[
  \sum_{k=1}^{n+1} k^2 = \frac{n(n+1)(2n+1)}{6} + (n+1)^2
  = \frac{(n+1)(n+2)(2n+3)}{6},
\]
which is exactly the claimed formula evaluated at $n+1$. By induction the
identity holds for all $n \geq 1$. \hfill\qed

\qpart
Use the formula above to evaluate $\sum_{k=1}^{10} k^2$.

Substituting $n=10$ gives
$\frac{10 \cdot 11 \cdot 21}{6} = 385$.

\end{question}

\begin{question}
Solve the second-order linear differential equation
\[
  \deriv{^2y}{x^2} - 5\deriv{y}{x} + 6y = 0.
\]

The auxiliary equation is $\lambda^2 - 5\lambda + 6 = 0$, which factorises
to $(\lambda - 2)(\lambda - 3) = 0$, giving roots $\lambda = 2$ and
$\lambda = 3$. Since these are real and distinct, the general solution is
\[
  y = A\e^{2x} + B\e^{3x}.
\]

\end{question}

\begin{question}[4]
\qsubpart
Find the eigenvalues of the matrix
$\begin{pmatrix} 4 & 1 \\ 2 & 3 \end{pmatrix}$.

The characteristic equation is
$(4-\lambda)(3-\lambda) - 2 = 0$, i.e.\ $\lambda^2 - 7\lambda + 10 = 0$,
which factorises to $(\lambda-2)(\lambda-5)=0$. The eigenvalues are
$\lambda = 2$ and $\lambda = 5$.

\qsubpart
Find an eigenvector corresponding to the smaller eigenvalue.

For $\lambda = 2$: $(4-2)v_1 + v_2 = 0 \Rightarrow v_2 = -2v_1$, so an
eigenvector is $\begin{pmatrix} 1 \\ -2 \end{pmatrix}$.

\end{question}

\end{document}
